\documentclass[11pt,a4paper]{article}

% ---------------------------------------------------------
% PREAMBLE (stable, warning-free, Overleaf-safe)
% ---------------------------------------------------------
\usepackage[utf8]{inputenc}
\usepackage[T1]{fontenc}
\usepackage{lmodern}
\usepackage{amsmath, amssymb, amsfonts}
\usepackage{bm}
\usepackage{geometry}
\usepackage{graphicx}
\usepackage{hyperref}
\usepackage{cite}
\usepackage{authblk}
\usepackage{physics}
\usepackage{mathtools}

\geometry{margin=2.4cm}

\title{\textbf{Companion L: The Growth Index in the Relaxation-Driven Cyclic Cosmology}}
\author{Michael Lehmann \\ \texttt{mi.lehmann@gmx.de}}
\date{April 2026}

\begin{document}
\maketitle

\begin{abstract}
The growth index $\gamma$ provides a compact description of the dynamical evolution 
of cosmic structure. In the standard cosmological model, $\gamma$ is nearly 
constant, reflecting the scale-independent growth of density perturbations. In the 
Relaxation-Driven Cyclic Cosmology (RDCC), the scale-dependent evolution of the 
gravitational potential modifies the growth rate in a way that renders the growth 
index scale-dependent and time-dependent. This Companion develops a continuous 
description of the growth index in RDCC, deriving how the relaxation scale 
influences the growth rate, the scale dependence of $\gamma$, and the observational 
signatures in redshift-space distortions, weak lensing, and velocity fields. The 
resulting framework provides a dynamical fingerprint of the RDCC gravitational 
field.
\end{abstract}
% ---------------------------------------------------------
\section{Introduction}

The growth index $\gamma$ is defined through the relation
\begin{equation}
    f(a) = \Omega_{m}(a)^{\gamma},
\end{equation}
where $f(a)$ is the linear growth rate. In the standard cosmological model, the 
growth index is nearly constant, with $\gamma \approx 0.55$ for a Universe 
dominated by cold dark matter and a cosmological constant.

In RDCC, the gravitational potential evolves differently. The relaxation mechanism 
suppresses the decay of small-scale modes and enhances the coherence of large-scale 
modes. This modifies the growth rate in a scale-dependent manner, rendering the 
growth index both scale-dependent and time-dependent. The growth index becomes a 
direct probe of the relaxation scale $k_{\mathrm{rel}}$ and the temporal structure 
of the RDCC gravitational field.
% ---------------------------------------------------------
\section{The RDCC Growth Rate}

The linear growth rate is defined as
\begin{equation}
    f(k,a)
    = \frac{d\ln D(k,a)}{d\ln a},
\end{equation}
where $D(k,a)$ is the growth factor.

In RDCC, the growth factor evolves as
\begin{equation}
    D_{\mathrm{RDCC}}(k,a)
    = D(a)
      \left[
        1 - \chi_{D}
        \left(\frac{k}{k_{\mathrm{rel}}}\right)^{\alpha}
      \right].
\end{equation}

The growth rate becomes
\begin{equation}
    f_{\mathrm{RDCC}}(k,a)
    = f(a)
      \left[
        1 - \chi_{f}
        \left(\frac{k}{k_{\mathrm{rel}}}\right)^{\alpha}
      \right].
\end{equation}

This modification introduces a scale dependence that reflects the relaxation 
mechanism.
% ---------------------------------------------------------
\section{The RDCC Growth Index}

The growth index is defined implicitly by
\begin{equation}
    f_{\mathrm{RDCC}}(k,a)
    = \Omega_{m}(a)^{\gamma_{\mathrm{RDCC}}(k,a)}.
\end{equation}

Solving for $\gamma_{\mathrm{RDCC}}$ yields
\begin{equation}
    \gamma_{\mathrm{RDCC}}(k,a)
    = \frac{
        \ln f(a)
        + \ln\left[
            1 - \chi_{f}
            \left(\frac{k}{k_{\mathrm{rel}}}\right)^{\alpha}
        \right]
    }{
        \ln \Omega_{m}(a)
    }.
\end{equation}

This expression shows that the growth index becomes scale-dependent and acquires a 
distinctive RDCC signature near the relaxation scale.
% ---------------------------------------------------------
\section{Scale Dependence and RDCC Signatures}

In the standard cosmological model, the growth index is nearly constant. In RDCC, 
the scale dependence becomes
\begin{equation}
    \frac{\partial \gamma_{\mathrm{RDCC}}}{\partial k}
    \propto
    -\chi_{f}\,\alpha\,
    \frac{k^{\alpha - 1}}{k_{\mathrm{rel}}^{\alpha}}.
\end{equation}

This modification produces:

- a suppressed growth index at small scales,  
- an enhanced growth index at large scales,  
- a characteristic turnover near $k_{\mathrm{rel}}$.

These signatures provide a direct observational probe of the relaxation scale.
% ---------------------------------------------------------
\section{Observational Probes}

The growth index can be measured through several observables:

\begin{itemize}
    \item \textbf{Redshift-space distortions (RSD)} measure $f\sigma_{8}$.
    \item \textbf{Weak lensing} measures the combination $f\Omega_{m}$.
    \item \textbf{Velocity fields} measure the velocity divergence.
    \item \textbf{Cluster abundance} depends on the growth history.
\end{itemize}

In RDCC, all of these observables acquire a scale-dependent modification that 
reflects the relaxation mechanism. The growth index therefore provides a 
dynamical fingerprint of the RDCC gravitational field.
% ---------------------------------------------------------
\section{Conclusion}

The growth index in the Relaxation-Driven Cyclic Cosmology reflects the 
scale-dependent evolution of the gravitational potential. The relaxation mechanism 
modifies the growth rate in a scale-dependent manner, rendering the growth index 
both scale-dependent and time-dependent. These effects produce distinctive 
signatures in redshift-space distortions, weak lensing, and velocity fields. The 
present Companion establishes the theoretical foundation for future studies of 
cosmic growth, dynamical evolution, and the interplay between matter and gravity in 
RDCC.

\bibliographystyle{apsrev4-2}
\nocite{*}
\bibliography{references}

\end{document}
